problem:  
Let \((\mathbb{H}^{n+1}(c), g_{\mathbb{H}})\) be the hyperbolic space of constant sectional curvature \(c < 0\), and let \(f:\mathbb{H}^{n+1}(c)\to\mathbb{R}\) be the height function with respect to a chosen point at infinity (so that \(\nabla f\) is the unit horospherical vector field and \(e^{-f}\) is the Busemann weight).  

A hypersurface \(M^n\subset\mathbb{H}^{n+1}(c)\) is called **\(f\)-minimal** if its weighted mean curvature satisfies
\[
H_f = H - \langle \nabla f, \nu \rangle = 0,
\]
where \(H\) is the mean curvature and \(\nu\) is the unit normal to \(M\). These are natural generalizations of self-shrinkers for mean curvature flow to negatively curved settings.  

Consider the class of compact \(f\)-minimal hypersurfaces in \(\mathbb{H}^{n+1}(c)\) that have **constant weighted scalar curvature**
\[
R_f := R + 2\Delta f - |\nabla f|^2 = \text{constant}.
\]
  
**Question:**  
Do there exist compact, non‑umbilical hypersurfaces satisfying these equations, or does the constancy of \(R_f\) force such an \(M\) to be totally umbilical in the metric \(g_{\mathbb{H}}\)?  
Equivalently, is the weighted scalar‑curvature condition sufficient to rigidify compact \(f\)-minimal hypersurfaces in hyperbolic space to geodesic spheres (up to isometries)?  
If not, can one describe a nontrivial family of examples (perhaps of rotational or warped type) realizing new compact \(f\)-minimal constant \(R_f\) solutions?

Why is it a "good" problem:  
This reformulation corrects the geometric inconsistencies while preserving the deep idea: searching for classification and rigidity of weighted minimal hypersurfaces in hyperbolic space, in direct analogy with the role played by Clifford tori in the sphere. The problem’s statement is conceptually simple—does constancy of a natural weighted curvature quantity imply geometric rigidity?—yet its analytic and geometric complexity is immense.  

It acts as a bridge between **weighted geometry** (Bakry–Émery theory), **minimal submanifold theory in space forms**, and **mean curvature flow solitons**, connecting static curvature functionals with dynamical soliton concepts. It also probes a profound frontier: whether hyperbolic geometry, when coupled with a natural Busemann weight, admits compact weighted‑minimal structures at all, and if so, whether they must mimic spherical models.  

Resolving it—either proving rigidity or constructing counterexamples—would unveil deep interactions between extrinsic curvature, weighting, and hyperbolic structure, marking a significant step in understanding weighted geometric variational problems beyond the Euclidean and spherical settings.